%% =====================================================================
%%  Solutions du chapitre 3 — Variables
%% =====================================================================
\corrigesChapitre{3}{Variables}

\begin{corrigeatoi}{1}
\begin{suivi}{2}
\otTH{Line} & \otTH{Instruction} & \otTH{x} & \otTH{y}\\\midrule
1 & \py{x = 3} & \nv{3} & \\
2 & \py{y = x + 4} & \av{3} & \nv{7}\\
3 & \py{x = y * 2} & \nv{14} & \av{7}\\
4 & \py{y = x - y} & \av{14} & \nv{7}\\
\end{suivi}
\begin{code}{file=trace2.py}
x = 3
y = x + 4
x = y * 2
y = x - y
print(x, y)
\end{code}
\begin{sortie}
14 7
\end{sortie}
At line~4, \py{y} keeps the value 7, but it is a coincidence: it is calculated again from the values
of that moment, $14 - 7$.
\end{corrigeatoi}

\begin{corrigeatoi}{2}
\begin{code}{file=points_game.py}
points = 0
points += 50
points -= 20
points *= 3
print("Final score:", points)
\end{code}
\begin{sortie}
Final score: 90
\end{sortie}
\end{corrigeatoi}

\begin{corrigeatoi}{3}
One more night also means one more day of meals: change \py{nights} \emph{and} \py{days}
(lines~5 and~6).
\begin{code}{file=holiday_planner.py, marks={4-6}}
# holiday_planner.py - the Lisbon budget, now with variables

# Data
travellers = 3
nights = 7
days = 8
flight_price = 89.5      # return ticket, per person
room_price = 74          # per night, for the room
food_per_day = 25        # per person
museum_pass = 21.5       # per person
months = 6               # months left to save

# Calculations
flights = travellers * flight_price
hotel = nights * room_price
food = travellers * days * food_per_day
museums = travellers * museum_pass
total = flights + hotel + food + museums

# Display
print("=== HOLIDAY IN LISBON ===")
print("Travellers:", travellers)
print("Flights:", flights)
print("Hotel:", hotel)
print("Food:", food)
print("Museums:", museums)
print("Total:", total)
print("Per person:", round(total / travellers, 2))
print("To save each month:", round(total / months, 2))
\end{code}
\begin{sortie}
=== HOLIDAY IN LISBON ===
Travellers: 3
Flights: 268.5
Hotel: 518
Food: 600
Museums: 64.5
Total: 1451.0
Per person: 483.67
To save each month: 241.83
\end{sortie}
\piege{changing only \py{nights}: the hotel would be right, but the food would still be counted for
7 days.}
\end{corrigeatoi}

\begin{corrige}{1}
\begin{suivi}{2}
\otTH{Line} & \otTH{Instruction} & \otTH{a} & \otTH{b}\\\midrule
1 & \py{a = 4} & \nv{4} & \\
2 & \py{b = a * 3} & \av{4} & \nv{12}\\
3 & \py{a = b - a} & \nv{8} & \av{12}\\
4 & \py{b = b // a} & \av{8} & \nv{1}\\
\end{suivi}
\begin{code}{file=trace3.py}
a = 4
b = a * 3
a = b - a
b = b // a
print(a, b)
\end{code}
\begin{sortie}
8 1
\end{sortie}
\retiens{at each line, use the values of the line \emph{above} in the trace table: they are the values
the variables have at that moment.}
\end{corrige}

\begin{corrige}{2}
\begin{code}{file=follow.py}
x = 10
y = x
x = x + 5
print(x, y)
\end{code}
\begin{sortie}
15 10
\end{sortie}
\pourquoi{at line~2, \py{y} is attached to the object \py{10}. At line~3, \py{x} is attached to a new
object, \py{15}; this does not move the label \py{y}, which still refers to \py{10}.}
\end{corrige}

\begin{corrige}{3}
\begin{reponses}
\item \py{age}: valid and well chosen.
\item \py{2fast}: \rep{invalid}, a name cannot start with a digit.
\item \py{first\_name}: valid, perfect snake\_case.
\item \py{first-name}: \rep{invalid}, Python reads a subtraction.
\item \py{\_count}: valid. By convention, a name that starts with an underscore is for internal use;
  you will meet this in volume~3.
\item \py{class}: \rep{invalid}, it is a keyword.
\item \py{Class}: valid, since Python is case-sensitive, but names with a capital letter are reserved
  by convention for \emph{classes} (volume~3).
\item \py{my score}: \rep{invalid}, spaces are forbidden.
\item \py{totalPrice}: valid, but not the Python style: write \py{total\_price}.
\item \py{MAX\_SPEED}: valid, the convention for a constant.
\end{reponses}
\end{corrige}

\begin{corrige}{4}
\begin{code}{file=rectangle.py, cible=rectangle}
width = 7.5
height = 4
area = width * height
perimeter = 2 * (width + height)
print("Area:", area, "square metres")
print("Perimeter:", perimeter, "metres")
\end{code}
\pourquoi{the results are floats (\py{30.0}, \py{23.0}) because \py{width} is a float.}
\piege{writing \py{2 * width + height}: without parentheses, only the width is doubled.}
\end{corrige}

\begin{corrige}{5}
\begin{code}{file=shopping.py, cible=courses}
bread_price = 1.25
milk_price = 0.75
milk_bottles = 3
total = bread_price + milk_price * milk_bottles
print("Total:", total)
\end{code}
\pourquoi{the multiplication comes first: $1.25 + 0.75 \times 3 = 1.25 + 2.25 = 3.5$. In the last
line, \py{total} is written without quotes: we want its value, not the word \enquote{total}.}
\end{corrige}

\begin{corrige}{6}
\demarche{read the message: \py{hour} does not exist, Python suggests \py{hours}. Correct it and run
again: a second \py{NameError} appears on line~4, because of the capital letter of
\py{Weekly\_pay}.}
\begin{code}{file=weekly_pay.py, cible=salaire}
hourly_rate = 12.5
hours = 35
weekly_pay = hourly_rate * hours
print("Weekly pay:", weekly_pay)
\end{code}
\retiens{a name must be written exactly the same way, letter for letter and with the same case,
every time you use it.}
\end{corrige}

\begin{corrige}{7}
\begin{code}{file=swap.py, cible=echange}
left = "red"
right = "blue"
print("Before:", left, right)

# first method: with a temporary variable
temp = left
left = right
right = temp
print("After the first swap:", left, right)

# second method: Python's multiple assignment
left, right = right, left
print("After the second swap:", left, right)
\end{code}
\pourquoi{\py{temp} keeps the old value of \py{left} before it is replaced. The second method swaps
again, which brings back the initial order.}
\piege{writing \py{left = right} then \py{right = left}: after the first line, both names refer to
\py{"blue"}, and \py{"red"} is lost.}
\end{corrige}

\begin{corrige}{8}
\begin{code}{file=speed.py, cible=vitesse}
distance = 150     # in km
duration = 2.5     # in hours
speed = distance / duration
print("Average speed:", speed, "km/h")
print("That is", round(speed / 3.6, 2), "m/s")
\end{code}
\end{corrige}

\begin{corrige}{9}
\begin{suivi}{1}
\otTH{Line} & \otTH{Instruction} & \otTH{points}\\\midrule
1 & \py{points = 10} & \nv{10}\\
2 & \py{points += 5} & \nv{15}\\
3 & \py{points *= 2} & \nv{30}\\
4 & \py{points -= 4} & \nv{26}\\
5 & \py{points //= 3} & \nv{8}\\
\end{suivi}
\begin{code}{file=points.py}
points = 10
points += 5
points *= 2
points -= 4
points //= 3
print(points)
\end{code}
\begin{sortie}
8
\end{sortie}
\end{corrige}

\begin{corrige}{10}
\demarche{the recipe is for 4 people; for \py{guests} people, every quantity is multiplied by
\py{guests / 4}.}
\begin{code}{file=pancakes.py, cible=crepes}
guests = 6
factor = guests / 4
print("Pancakes for", guests, "people:")
print("Flour:", 250 * factor, "g")
print("Milk:", 500 * factor, "ml")
print("Eggs:", 2 * factor)
\end{code}
For 2 guests, only the first line changes:
\begin{code}{file=pancakes.py, variante, first=1, marks={1}}
guests = 2
\end{code}
\begin{sortie}
Pancakes for 2 people:
Flour: 125.0 g
Milk: 250.0 ml
Eggs: 1.0
\end{sortie}
\retiens{thanks to the variable \py{guests}, one single change adapts the whole recipe.}
\end{corrige}

\begin{corrige}{11}
\begin{memoire}[after each line]
\node[font=\otSemi\fontsize{8}{9}\selectfont,text=otViolet] at (1.6,1.5) {after line 1};
\nomvar{a}{(0,0.5)}{a1}
\objet{int}{3}{(3,0.5)}{o1}
\lien{a1}{o1}
\node[font=\otSemi\fontsize{8}{9}\selectfont,text=otViolet] at (6.9,1.5) {after line 2};
\nomvar{a}{(5.3,0.5)}{a2}
\nomvar{b}{(5.3,-0.8)}{b2}
\objet{int}{3}{(8.3,-0.15)}{o2}
\lien{a2}{o2}
\lien{b2}{o2}
\node[font=\otSemi\fontsize{8}{9}\selectfont,text=otViolet] at (12.2,1.5) {after line 3};
\nomvar{a}{(10.6,0.5)}{a3}
\nomvar{b}{(10.6,-0.8)}{b3}
\objet{int}{7}{(13.6,0.5)}{o3}
\objet{int}{3}{(13.6,-0.8)}{o4}
\lien{a3}{o3}
\lien{b3}{o4}
\end{memoire}
\begin{code}{file=memory.py}
a = 3
b = a
a = 7
print(a, b)
\end{code}
\begin{sortie}
7 3
\end{sortie}
\pourquoi{line~2 attaches a second name, \py{b}, to the object \py{3}. Line~3 moves the label \py{a}
onto a new object, \py{7}; \py{b} still refers to \py{3}.}
\end{corrige}

\begin{corrige}{12}
\demarche{\py{n \% 10} gives the last digit; \py{n // 10} removes it. Combine them to reach each
digit: \py{n // 10 \% 10} is the tens digit, \py{n // 100 \% 10} the hundreds digit.}
\begin{code}{file=digits.py, cible=chiffres}
n = 4728
thousands = n // 1000
hundreds = n // 100 % 10
tens = n // 10 % 10
units = n % 10
print("Sum of the digits:", thousands + hundreds + tens + units)
print("Backwards:", units * 1000 + tens * 100 + hundreds * 10 + thousands)
\end{code}
\pourquoi{for 4728: \py{thousands} is 4, \py{hundreds} is 7, \py{tens} is 2 and \py{units} is 8.
The sum is $4 + 7 + 2 + 8 = 21$. To write the number backwards, the units become the thousands, and
so on.}
\retiens{name your intermediate results: the last two lines read almost like the statement of the
exercise.}
\end{corrige}

\begin{corrigequiz}
\begin{reponses}[5]
\item b
\item a
\item c
\item b
\item b
\item b
\item c
\item b
\item b
\item c
\end{reponses}
\pourquoi{2: an assignment displays nothing in the Shell. 5: \py{b} keeps the value 5 when \py{a}
changes. 8: the quotes make \py{"42"} a piece of text, not a number. 9: a variable only exists once
its line has run.}
\end{corrigequiz}
